\documentclass[11pt]{article}
\usepackage[T1]{fontenc}
\usepackage{amsmath,amssymb,amsthm}
\usepackage{graphicx}
\usepackage[margin=1in]{geometry}
\usepackage{microtype}
\usepackage[hidelinks]{hyperref}

\newtheorem{theorem}{Theorem}
\newtheorem{corollary}{Corollary}
\newtheorem{proposition}{Proposition}
\newtheorem{lemma}{Lemma}
\newtheorem{definition}{Definition}
\newcommand{\R}{\mathbb R}
\newcommand{\SX}{S_X}
\newcommand{\BX}{B_X}
\newcommand{\perpI}{\mathrel{\perp_I}}
\newcommand{\perpP}{\mathrel{\perp_P}}
\newcommand{\Lcal}{\mathcal L}

\title{A universal orthogonality-transfer principle\\
for geometric constants}
\author{Automated mathematics run \texttt{fa\_banach\_001}, lane 10}
\date{17 August 2026}

\begin{document}
\maketitle

\begin{center}
\fbox{\parbox{0.92\textwidth}{
\textbf{Status: candidate full structural answer, likely valid; novelty not
certified.}  Under the literal meaning of ``equivalently expressed'', all
finite-point unit-sphere constants admit exact isosceles-orthogonal
coordinates.  This gives arbitrary many-parameter extensions of
$L'_{\mathrm{YJ}}$, a necessary-and-sufficient criterion for other
orthogonalities, a new Pythagorean-orthogonality transfer, and exact operator
versions.  The source questions are informal, so claims about which formulas
are \emph{natural} or \emph{useful} remain outside the result.
}}
\end{center}

\section{Current source questions and version boundary}

Wang, Liu, and Bao \cite{WangLiuBao} ask in Problems~5.8--5.12 which geometric
constants have orthogonality-based representations, whether their
two-parameter $L'_{\mathrm{YJ}}$ construction has more-parameter analogues,
whether unit-ball or unit-sphere restrictions can generally be encoded by
isosceles orthogonality, whether other orthogonalities work, and how operator
versions should be treated.

The current source is arXiv:2507.17122v5, revised 18 July 2026.  This matters:
v5 itself adds a Birkhoff--James rectification theorem and exact formulas for
several constants.  Those results are prior source material here, not claimed
as new.  The five questions below remain on printed page~22 of v5.

\begin{figure}[ht]
 \centering
 \includegraphics[width=\textwidth]{figures/open_problems_crop.png}
 \caption{Problems~5.8--5.12 in arXiv:2507.17122v5, printed page~22.}
\end{figure}

\section{Proof intuition}

There is an exact, norm-independent change of coordinates behind every
isosceles formula in the source.  Given two unit vectors $p,q$, their half-sum
$u=(p+q)/2$ and half-difference $v=(p-q)/2$ satisfy
$\|u+v\|=\|u-v\|=1$; hence $u\perpI v$.  Conversely $p=u+v$ and $q=u-v$.
This is a bijection, not an inequality.  Therefore every function, constraint,
parameter, supremum, or infimum on two sphere points transfers exactly.

For another orthogonality relation, the corresponding issue is precisely
surjectivity: can every vector $z$ be written as a longitudinal part
$\beta x$ plus a part $y$ orthogonal to $x$?  This ``rectifiability'' is both
necessary and sufficient for universal transfer.  Birkhoff--James
rectifiability is proved in v5 of the source.  Pythagorean orthogonality is a
new example: a one-dimensional convex forward-difference argument always
chooses the required $\beta$.

\section{The Hadamard transfer theorem}

Let $X$ be a nonzero real normed space and define the normalized isosceles
cone
\[
 \mathcal I_X:=\{(u,v)\in X^2:u\perpI v,\ \|u+v\|=1\}.
\]
The orthogonality condition automatically also gives $\|u-v\|=1$.

\begin{theorem}[Universal two-point transfer]\label{thm:hadamard}
The map
\[
 H:\SX^2\longrightarrow\mathcal I_X,
 \qquad H(p,q)=\left(\frac{p+q}{2},\frac{p-q}{2}\right)
\]
is a bijection with inverse $H^{-1}(u,v)=(u+v,u-v)$.  Consequently, for every
set of parameters $\Theta$ and every extended-real function
$F:\Theta\times\SX^2\to[-\infty,\infty]$,
\begin{align*}
 \sup_{\theta\in\Theta,\ p,q\in\SX}F(\theta,p,q)
 &=\sup_{\theta\in\Theta,\ (u,v)\in\mathcal I_X}
   F(\theta,u+v,u-v),\\
 \inf_{\theta\in\Theta,\ p,q\in\SX}F(\theta,p,q)
 &=\inf_{\theta\in\Theta,\ (u,v)\in\mathcal I_X}
   F(\theta,u+v,u-v).
\end{align*}
The same is true after imposing any predicate on the original variables and
its pullback under $H$.
\end{theorem}

\begin{proof}
If $p,q\in\SX$ and $(u,v)=H(p,q)$, then $u+v=p$ and $u-v=q$, so
$\|u+v\|=\|u-v\|=1$.  Thus $(u,v)\in\mathcal I_X$.  Conversely, if
$(u,v)\in\mathcal I_X$, then $p=u+v$ and $q=u-v$ belong to $\SX$, and the two
displayed linear maps undo each other.  Every assertion about functions,
extrema, parameters, and predicates follows by reindexing along this
bijection.
\end{proof}

\begin{corollary}[Finite sphere and ball constants]\label{cor:finite}
Every functional of finitely many unit vectors has an exact representation
using finitely many isosceles-orthogonal pairs, by applying
Theorem~\ref{thm:hadamard} to successive pairs of variables (and adding one
ignored dummy sphere variable if their number is odd).

Moreover, for every extended-real $G:\BX^2\to[-\infty,\infty]$,
\[
 \sup_{x,y\in\BX}G(x,y)
 =\sup_{\substack{(u,v)\in\mathcal I_X\\r,s\in[0,1]}}
 G\bigl(r(u+v),s(u-v)\bigr),
\]
and similarly for infima and constraints.
\end{corollary}

\begin{proof}
Only the ball formula needs comment.  Write a nonzero $x\in\BX$ as
$x=r p$, where $r=\|x\|\in(0,1]$ and $p\in\SX$; if $x=0$, take $r=0$ and
any $p\in\SX$.  Do the same for $y=sq$, and apply
Theorem~\ref{thm:hadamard} to $(p,q)$.  The converse inclusion is immediate.
\end{proof}

Thus, in the exact reparameterization sense, Problems~5.8 and 5.10 have the
maximal answer: all finite-point sphere constants transfer, and all
finite-point ball constants transfer after retaining their radial variables.

\section{Arbitrarily many parameters}

For an integer $N\geq1$ and a nonzero matrix
$A=(a_{\ell j})\in\R^{N\times2}$, define
\begin{equation}\label{eq:matrixconstant}
 \Lcal_A(X):=\frac{1}{\|A\|_{F}^{2}}
 \sup_{p,q\in\SX}\sum_{\ell=1}^{N}
 \|a_{\ell1}p+a_{\ell2}q\|^2.
\end{equation}
Up to common scaling, this is a family with $2N-1$ scalar parameters.

\begin{theorem}[Matrix-coefficient $L_{\mathrm{YJ}}$ family]\label{thm:matrix}
For every nonzero $A\in\R^{N\times2}$,
\begin{equation}\label{eq:matrixiso}
 \Lcal_A(X)=\frac{1}{\|A\|_F^2}
 \sup_{\substack{u\perpI v\\(u,v)\ne(0,0)}}
 \frac{\displaystyle\sum_{\ell=1}^{N}
 \|(a_{\ell1}+a_{\ell2})u+(a_{\ell1}-a_{\ell2})v\|^2}
 {\|u+v\|^2}.
\end{equation}
For
\[
 A=\begin{pmatrix}\tau&\upsilon\\ \upsilon&-\tau\end{pmatrix},
\]
the constant in \eqref{eq:matrixconstant} is exactly
$L'_{\mathrm{YJ}}(\tau,\upsilon,X)$ from the source.
\end{theorem}

\begin{proof}
For a nonzero isosceles pair put $r=\|u+v\|=\|u-v\|>0$ and
$p=(u+v)/r$, $q=(u-v)/r$.  Then
\[
 a_{\ell1}p+a_{\ell2}q
 =\frac{(a_{\ell1}+a_{\ell2})u+(a_{\ell1}-a_{\ell2})v}{r}.
\]
Theorem~\ref{thm:hadamard} gives \eqref{eq:matrixiso}.  For the displayed
two-row matrix, $\|A\|_F^2=2(\tau^2+\upsilon^2)$, and the numerator in
\eqref{eq:matrixconstant} is
$\|\tau p+\upsilon q\|^2+\|\upsilon p-\tau q\|^2$, which is the defining
formula for $L'_{\mathrm{YJ}}$.
\end{proof}

Since $N$ is arbitrary, Theorem~\ref{thm:matrix} answers Problem~5.9 with a
single family containing arbitrarily many independent coefficients.

\section{Which other orthogonalities transfer?}

Let $\mathrel{\perp_R}$ be any relation between a nonzero first vector and a
second vector in $X$.

\begin{definition}
The relation $\perp_R$ is \emph{rectifiable} if, for every $x\ne0$ and every
$z\in X$, there are $\beta\in\R$ and $y\in X$ such that
\[
 z=\beta x+y,\qquad x\perp_R y.
\]
It is \emph{sphere-rectifiable} if the same property holds whenever
$x=(p+q)/2\ne0$ and $z=(p-q)/2$ for $p,q\in\SX$.
\end{definition}

\begin{theorem}[Rectifiability characterization]\label{thm:rectification}
The following are equivalent for a relation $\perp_R$.
\begin{enumerate}
 \item The relation is rectifiable.
 \item The coordinate map
 \[
 (x,y,\beta)\longmapsto(x,\beta x+y),\qquad x\perp_R y,
 \]
 is onto $(X\setminus\{0\})\times X$.
 \item For every bounded function $\Phi:(X\setminus\{0\})\times X\to\R$,
 \[
 \sup_{x\ne0,z\in X}\Phi(x,z)
 =\sup_{\substack{x\perp_R y\\x\ne0,\ \beta\in\R}}
 \Phi(x,\beta x+y),
 \]
 and likewise for infima.
\end{enumerate}
The analogous universal transfer on $\SX^2$ holds if and only if the relation
is sphere-rectifiable, with antipodal pairs $(p,-p)$ included as a separate
boundary component.  Every rectifiable relation is sphere-rectifiable.
\end{theorem}

\begin{proof}
The first two assertions are the same surjectivity statement.  Surjectivity
implies the extrema identities by substitution.  Conversely, if a pair
$(x_0,z_0)$ is missed, take $\Phi$ to be its indicator: the left supremum is
$1$ and the right one is $0$.  The sphere statement applies the same argument
to
\[
 p=\frac{(1+\beta)x+y}{r},\qquad
 q=\frac{(1-\beta)x-y}{r},
\]
where the two numerator norms have the same positive value $r$.  A non-antipodal
pair has $x=(p+q)/2\ne0$ and is reached exactly when its half-difference can be
rectified.  Antipodal pairs have $x=0$ and form the stated boundary component.
\end{proof}

This is a necessary-and-sufficient formal answer to Problem~5.11.  The current
source proves that Birkhoff--James orthogonality is rectifiable by a norming
functional.  The next lemma supplies a different named example.

\begin{lemma}[Pythagorean orthogonality is rectifiable]\label{lem:pythagorean}
Define $x\perpP y$ by
\[
 \|x-y\|^2=\|x\|^2+\|y\|^2.
\]
Then $\perp_P$ is rectifiable on every real normed space.
\end{lemma}

\begin{proof}
Fix $x\ne0$ and $z\in X$, and put $q(t)=\|z-tx\|^2$.  The function $q$ is
continuous, convex, and satisfies
\[
 \frac{q(t)}{t^2}\longrightarrow\|x\|^2\qquad(|t|\to\infty).
\]
Its forward difference $d(t)=q(t+1)-q(t)$ is continuous and nondecreasing.
For $t>0$, convexity gives
\[
 d(t)\geq\frac{q(t)-q(0)}{t}\longrightarrow+\infty.
\]
For $t<-1$, comparison of the secant slopes on $[t,t+1]$ and $[t+1,0]$
gives
\[
 d(t)\leq\frac{q(0)-q(t+1)}{-t-1}\longrightarrow-\infty.
\]
Hence there is $\beta\in\R$ with $d(\beta)=\|x\|^2$.  Set
$y=z-\beta x$.  Since
\[
 \|y\|^2=q(\beta),\qquad \|x-y\|^2=q(\beta+1),
\]
we obtain $\|x-y\|^2=\|x\|^2+\|y\|^2$, as required.
\end{proof}

Theorems~\ref{thm:rectification} and Lemma~\ref{lem:pythagorean} now produce
exact Pythagorean-orthogonal representations for every two-vector full-space
constant, and—with the equal-radius endpoint condition—for every sphere
constant.

\section{Operator constants and moduli}

\begin{corollary}[Universal operator transfer]\label{cor:operator}
Let $T:X\to Y$ be a bounded linear operator, let $\Theta$ be any parameter
set, and let $\Psi:\Theta\times Y^2\to[-\infty,\infty]$.  Then
\[
 \sup_{\theta\in\Theta,\ p,q\in\SX}\Psi(\theta,Tp,Tq)
 =\sup_{\theta\in\Theta,\ (u,v)\in\mathcal I_X}
 \Psi\bigl(\theta,T(u+v),T(u-v)\bigr),
\]
with identical statements for infima and arbitrary constraints.  Pairwise
application gives the same result for every finite-point operator constant or
modulus, and Corollary~\ref{cor:finite} gives unit-ball versions.
\end{corollary}

\begin{proof}
Apply Theorem~\ref{thm:hadamard} to
$F(\theta,p,q)=\Psi(\theta,Tp,Tq)$.  Linearity gives the displayed form.
Products of the bijection handle more variables, and predicates pull back
exactly.
\end{proof}

For example, \eqref{eq:matrixconstant} has the operator analogue
\[
 \Lcal_A(T):=\frac1{\|A\|_F^2}\sup_{p,q\in\SX}
 \sum_{\ell=1}^N\|T(a_{\ell1}p+a_{\ell2}q)\|_Y^2,
\]
and Theorem~\ref{thm:matrix} remains valid after applying $T$ to every vector
in its numerator.  Constraints appearing in operator moduli of convexity or
smoothness transfer just as exactly.  This answers Problem~5.12 at the level
of every finite-configuration operator invariant.

\section{Verification, limitations, and novelty}

The proof uses only explicit bijections, surjectivity, and one-variable
convexity.  The main checks are: the zero half-sum occurs exactly at antipodal
sphere pairs and is retained as a boundary component; a nonzero isosceles pair
has $\|u+v\|>0$; the matrix normalization recovers the source denominator; and
the Pythagorean forward difference has the correct two infinite limits.  No
computational or unproved dependency remains.

The result gives a full answer only under the precise interpretation that an
``equivalent expression'' means an exact reparameterization, with auxiliary
scalar variables allowed.  The source does not define a narrower notion such
as formula complexity, naturality, or usefulness.  If one of those was
intended, the theorem is instead a complete structural baseline: it says
existence is automatic and identifies what additional criterion must be
specified for a nontrivial classification.

The run indexes and a bounded arXiv/web search through 17 August 2026 found
many particular orthogonal constants but no universal finite-point theorem,
arbitrary matrix family, rectifiability characterization, or Pythagorean
transfer.  Searches used the exact title and problems, $L'_{\mathrm{YJ}}$ with
matrix/multiple parameters, unit-sphere transfer, other orthogonalities, and
operator versions; close papers included arXiv:2111.08392, 2504.00826,
2509.23319, 2602.13974, 2606.01064, and 2606.01068.  Novelty is plausible, not
certified.

Human review should focus on whether the literal formalization matches the
authors' intended breadth, the necessity direction of
Theorem~\ref{thm:rectification}, and the secant-slope limits in
Lemma~\ref{lem:pythagorean}.  Subject to those checks and a specialist search,
promote as a full structural resolution of Problems~5.8--5.12.

\begingroup
\raggedright
\begin{thebibliography}{9}
\bibitem{WangLiuBao}
Q.~Liu, Y.~Wang, and M.~Bao,
\emph{How orthogonality influences geometric constants},
arXiv:2507.17122v5, Problems~5.8--5.12 (2025--2026).
\end{thebibliography}
\endgroup

\end{document}
